
\clearpage
\section{Eliminating the loan, retaining the currency exposure}
\label{sec:intragroup-fx-example}

An intragroup loan disappears from the consolidated balance sheet, yet its
exchange gain can remain in group profit. The reason becomes visible when
we follow both the loan and the money it financed. IAS~21.45 requires the
exchange result to survive elimination; paragraph~32 determines when it
belongs in other comprehensive income instead
\citep[IAS~21.32--33, .45]{eu20231803}.

Consider a parent whose functional and presentation currency is USD and a
wholly owned subsidiary whose functional currency is EUR. Control and both
functional currencies are established assumptions. At an opening rate of
USD1.20 per EUR, the parent's balances are cash 56, an investment in the
subsidiary carried at historical cost 24, and a EUR100 loan receivable worth
120. Its equity is therefore 200. The subsidiary holds EUR120 in cash,
owes EUR100 to the parent and has EUR20 of paid-in equity. At the opening
rate, its equity matches the parent's investment. Eliminating that
investment and the loan leaves group cash and equity of 200.

We isolate an exchange-rate remeasurement of these established balances,
holding every other measurement input fixed. This is neither the initial
valuation of a new loan nor a complete reporting period. Interest, tax,
impairment, hedges and new transactions are outside this calculation;
IAS~21.24 leaves the underlying carrying amount to the relevant measurement
standards. Thus EUR100 is a given monetary carrying amount, not an assumption
that an interest-free long-term loan must be recognised at its face value
\citep[IAS~21.23--24, .50]{eu20231803}. All following statement amounts are
in USD unless explicitly marked EUR.

Now let the rate rise to USD1.30 per EUR. The parent remeasures its receivable
from 120 to 130: debit the receivable 10 and credit exchange gain 10. The
subsidiary records no exchange gain on its EUR loan in its own EUR books.
For group presentation, its cash becomes $120\times1.30=156$ and its payable
becomes $100\times1.30=130$. Its historical capital remains 24. The remaining
2 is translation OCI, explained by its EUR20 net assets and the rate change:
$20\times(1.30-1.20)=2$
\citep[IAS~21.28, .33, .39--41]{eu20231803}.

Consolidation now debits the translated loan payable 130 and credits the
parent's receivable 130. It also debits the subsidiary's historical capital
24 and credits the parent's investment 24. These entries eliminate the
internal balances; neither entry removes the parent's exchange gain. Group
cash is $56+156=212$, with no external liability in this example. Opening
equity 200, exchange profit 10 and translation OCI 2 reconcile to the same
212. Deleting the profit merely because the receivable was internal would
leave equity short by 10. The group still owns the euros whose dollar value
has risen \citep[IFRS~10.B86; IAS~21.45]{eu20231803}.

\clearpage
The loan's settlement facts can change the presentation of that 10.
In the ordinary-loan case, settlement is planned in the foreseeable future,
so the gain stays in consolidated profit. In a separate conditional case,
assume settlement is neither planned nor likely in the foreseeable future
and the loan forms part of the net investment under IAS~21.15. This is a
substantive condition, not an election created by calling a loan long term;
trade receivables and payables are explicitly excluded. The parent's separate
gain is still 10. In the consolidated statements, debit exchange gain 10
and credit OCI 10, giving total OCI of 12 and no exchange profit. The two
cases in Table~\ref{tab:intragroup-fx} have identical group cash and total
equity but different profit and OCI
\citep[IAS~21.15, .32--33]{eu20231803}.

\begin{table}[!ht]
\centering
\caption{Settlement facts change where the exchange result appears.}
\label{tab:intragroup-fx}
\begin{tabular}{@{}lrr@{}}
\toprule
USD amount & Ordinary loan & Qualifying net-investment loan \\
\midrule
Parent's separate exchange gain & 10 & 10 \\
Consolidated exchange profit & 10 & 0 \\
Consolidated OCI & 2 & 12 \\
Consolidated closing cash and equity & 212 & 212 \\
\bottomrule
\end{tabular}
\end{table}

Neither case includes a cash receipt. The same EUR120 remains in the
subsidiary, so its translated value rises by
$120\times(1.30-1.20)=12$. IAS~7.28 presents this exchange effect separately
from operating, investing and financing cash flows. Group cash reconciles
as opening 200, transaction cash flows zero and exchange effect 12, giving
closing 212 \citep[IAS~7.25--28]{eu20231803}. Adding a further gain of 12
would count the same currency movement twice.

The representation must therefore retain both entities' functional currencies,
the monetary balance, the settlement assessment and the origin of each OCI
component. If the settlement facts are unknown, these conditional presentations
remain separate even when they agree on total equity. The group adjustment
does not overwrite the parent's separate journal, and the OCI history remains
available for the applicable later disposal treatment. An equal consolidated
total cannot establish that the histories answer the same profit, cash-flow
or disposal question.
